\section{泛函族的一致有界性}
\label{sec:lfunc-uniform-boundedness}

设$\{\ell_\alpha\}_{\alpha\in I}$是$X^*$中的一族泛函，并且在每个固定的向量上读数都有界：对每个$x$，$\sup_\alpha\lvert\ell_\alpha(x)\rvert<\infty$。这称为逐点有界。我们想知道能否推出$\sup_\alpha\lVert\ell_\alpha\rVert<\infty$，即一致有界。两者的量词次序不同：前者是“对每个$x$，存在常数$C_x$”，后者是“存在常数$C$，对每个单位向量$x$”。下一节研究弱收敛时，正需要从逐个读数的信息得到这样的统一控制。

\paragraph{不完备时可以失败}
在只有有限项非零的数列组成的空间$c_{00}$上配最大值范数，令$\ell_n(x)=nx_n$。对固定的$x$，$n$超过它最后一个非零坐标后$\ell_n(x)=0$，所以$\sup_n\lvert\ell_n(x)\rvert<\infty$；可是在$\bm e_n$上读数为$n$，$\lVert\ell_n\rVert=n$，没有统一上界。$c_{00}$不完备，例如$(1,1/2,\ldots,1/n,0,0,\ldots)$是Cauchy列，极限$(1/n)_n$不在$c_{00}$中。下面会看到，完备性正是把逐点有界转为一致有界的关键。

\paragraph{需要的工具}
第~\ref{chap:metric-spaces}章的Baire纲定理说：在非空完备度量空间中，可数个内部为空的闭集不能覆盖整个空间。换句话说，若非空完备度量空间写成可数个闭集之并，则至少有一个闭集含有某个开球。证明的思路是用嵌套闭球逐个避开各个闭集，完备性保证球心的极限存在；我们直接使用这一结论。

\begin{theorem}[泛函的一致有界原理]
\label{thm:lfunc-uniform-boundedness}
设$X$是Banach空间，$\{\ell_\alpha\}_{\alpha\in I}\subset X^*$，并且对每个$x\in X$有$\sup_\alpha\lvert\ell_\alpha(x)\rvert<\infty$。则$\sup_\alpha\lVert\ell_\alpha\rVert<\infty$。
\end{theorem}

这一结论称为\term{一致有界原理}（uniform boundedness principle）。证明按“闭集覆盖、找内点、平移到原点、缩放到单位球”的顺序进行。

\begin{proof}
\emph{闭集覆盖。}对正整数$k$令$F_k=\{x\in X:\sup_\alpha\lvert\ell_\alpha(x)\rvert\leq k\}=\bigcap_\alpha\{x:\lvert\ell_\alpha(x)\rvert\leq k\}$。每个$\ell_\alpha$连续，所以交集中的每个集合都是闭集，$F_k$也是闭集。逐点有界说明每个$x$属于某个$F_k$，于是$X=\bigcup_kF_k$。

\emph{找内点。}$X$完备且非空，由Baire纲定理，存在$k_0$和开球$B(x_0,r)\subset F_{k_0}$。

\emph{平移到原点。}对任意$\lVert z\rVert<r$，$x_0+z$与$x_0$都在$F_{k_0}$中，所以对每个$\alpha$，
\[
\lvert\ell_\alpha(z)\rvert=\lvert\ell_\alpha(x_0+z)-\ell_\alpha(x_0)\rvert\leq2k_0 .
\]

\emph{缩放到单位球。}对单位向量$x$及$0<s<r$，$\lVert sx\rVert<r$，由线性$\lvert\ell_\alpha(x)\rvert=s^{-1}\lvert\ell_\alpha(sx)\rvert\leq2k_0/s$。令$s\to r$，得$\lVert\ell_\alpha\rVert\leq2k_0/r$对一切$\alpha$成立。
\end{proof}

\paragraph{完备性用在哪里}
在$c_{00}$的例子中，$F_k$同样是覆盖空间的闭集，但没有一个含有开球：任何开球$B(x_0,r)$中都有$x_0+\frac r2\bm e_n$，取$n$足够大使它不属于$F_k$。论证失败的位置正是Baire纲定理，而Baire纲定理失败是因为空间不完备。同一个定理把线性、连续性和完备性组合起来，完成了量词的交换。

把$\ell_\alpha$换成取值在赋范空间中的有界线性算子，同样的四步论证给出算子族的一致有界原理；它与开映射定理一起，是第~\ref{chap:linear-operators}章研究完备空间中算子定理的出发点。本章只需要泛函的情形，下一节把它用于弱收敛列。
